\chapter{Code: Adversarial Mutators}
\label{apndx:code}

This appendix gives a before-and-after listing for each of the seven
semantics-preserving mutators described in Chapter~\ref{ch:impl}. Every
transformation leaves runtime behavior unchanged while breaking one detector
shortcut. The two core validation procedures, the Layer-1 sink-presence filter
and the Layer-3 diff filter, appear as Listings~\ref{lst:sinkfilter}
and~\ref{lst:difffilter} in Chapter~\ref{ch:impl}.

\section{Sink-Preserving Family (M1 to M4)}

These mutators alter the code around the sink but leave the sink identifier and
its arguments intact.

\begin{figure}[H]
\begin{verbatim}
# Before
def get_user(user_id):
    query = f"SELECT * FROM users WHERE id = {user_id}"
    return cursor.execute(query)

# After
def get_user(user_id):
    _unused_var = sum([])
    query = f"SELECT * FROM users WHERE id = {user_id}"
    if False:
        pass
    return cursor.execute(query)
\end{verbatim}
\caption{M1 dead-code injection. Two side-effect-free statements are added; the
return value cannot change. Attacks positional reasoning.}
\label{lst:m1}
\end{figure}

\begin{figure}[H]
\begin{verbatim}
# Before
return cursor.execute(f"SELECT * FROM users WHERE id = {uid}")

# After
return cursor.execute(
    "SEL" + "ECT * FROM" + " users WHERE id = " + f"{uid}"
)
\end{verbatim}
\caption{M2 string-literal splitting. The substring \texttt{"SELECT"} no longer
appears in the source, but the runtime SQL string is byte-identical. Attacks
substring matching. Held out as a test-time probe.}
\label{lst:m2}
\end{figure}

\begin{figure}[H]
\begin{verbatim}
# Before
def get_user(user_id):
    query = f"SELECT * FROM users WHERE id = {user_id}"
    return cursor.execute(query)

# After
def get_user(account_pk):
    sql = f"SELECT * FROM users WHERE id = {account_pk}"
    return cursor.execute(sql)
\end{verbatim}
\caption{M3 identifier rename. Locals and parameters are renamed to synonyms in
all their references. Attacks identifier memorization.}
\label{lst:m3}
\end{figure}

\begin{figure}[H]
\begin{verbatim}
# Before
def get_user(user_id):
    query = f"SELECT * FROM users WHERE id = {user_id}"
    return cursor.execute(query)

# After
def get_user(user_id):
    def _wrapped_execute_138(*_a, **_kw):
        return cursor.execute(*_a, **_kw)

    query = f"SELECT * FROM users WHERE id = {user_id}"
    return _wrapped_execute_138(query)
\end{verbatim}
\caption{M4 wrapper extraction. The sink call is moved into a one-line nested
helper that forwards arguments. Attacks single-function analysis.}
\label{lst:m4}
\end{figure}

\section{Sink-Targeted Family (M5, M6)}

These mutators rewrite the sink call itself, so the sink token survives only
inside a string literal. Both are held out as test-time probes.

\begin{figure}[H]
\begin{verbatim}
# Before
def run_cmd(user_arg):
    return os.system("ls " + user_arg)

# After
def run_cmd(user_arg):
    return getattr(os, "system")("ls " + user_arg)
\end{verbatim}
\caption{M5 sink attribute obfuscation. An attribute sink
\texttt{value.attr(args)} becomes \texttt{getattr(value, "attr")(args)}. Attacks
token-level sink matching for CWE-78, 89, 918, and 502.}
\label{lst:m5}
\end{figure}

\begin{figure}[H]
\begin{verbatim}
# Before
def render_template(spec):
    return eval(spec)

# After
def render_template(spec):
    return __import__("builtins").__dict__["eval"](spec)
\end{verbatim}
\caption{M6 sink via globals. A builtin sink \texttt{eval(arg)} is rewritten
through a \texttt{builtins} dictionary lookup, as shown above, so the sink token
survives only as a string literal. Attacks the same token-level match on the
CWE-94 bare-builtin position.}
\label{lst:m6}
\end{figure}

\section{Dataflow Family (M7)}

This mutator leaves the sink and its argument value unchanged and instead breaks
the static lineage from source to sink.

\begin{figure}[H]
\begin{verbatim}
# Before
def get_user(user_id):
    return cursor.execute(
        "SELECT * FROM users WHERE id = " + user_id
    )

# After
def get_user(user_id):
    return cursor.execute(
        {"_a": "SELECT * FROM users WHERE id = " + user_id}["_a"]
    )
\end{verbatim}
\caption{M7 taint-through-dict. The sink's first argument is routed through a
one-entry dict that is constructed and immediately re-extracted, so the value is
bitwise identical. Attacks taint tracking that does not propagate through dict
reads. Held out as a test-time probe.}
\label{lst:m7}
\end{figure}
